\chapter{Título según el objeto de estudio}

Cada capítulo debe iniciar con uno o dos párrafos de texto llano. En ellos se presenta brevemente el contenido, se indica el objetivo del capítulo y se anticipan los aspectos que serán desarrollados.

Este segundo párrafo puede utilizarse para delimitar el alcance del capítulo y conectar los conceptos que se desarrollarán en las secciones siguientes.

\section{Concepto 1}

Aquí se desarrolla el concepto principal. Para una cita narrativa puedes escribir: \textcite{ejemplo2020} sostiene que el concepto debe entenderse dentro de su contexto teórico.

También puedes usar una cita parentética al final de una idea \parencite{ejemplo2020}.

Para una cita textual corta: ``texto citado'' \parencite[87]{ejemplo2020}.

\section{Concepto 2}

Desarrolla los conceptos necesarios para comprender el objeto de estudio. No dejes subtítulos sin texto explicativo.

\section{Concepto 3}

Antes de una tabla, debes mencionarla desde el texto. Los datos de ejemplo se presentan en la Tabla~\ref{tab:ejemplo}.

\begin{table}[htbp]
  \caption{Ejemplo de presentación de datos en una tabla APA}
  \label{tab:ejemplo}
  \begin{spacing}{1.0}
  \centering
  \begin{tabular}{lccc}
    \toprule
    Categoría & Grupo A & Grupo B & Total \\
    \midrule
    Caso 1 & 10 & 12 & 22 \\
    Caso 2 & 8  & 9  & 17 \\
    \midrule
    Total  & 18 & 21 & 39 \\
    \bottomrule
  \end{tabular}
  \end{spacing}
  \notaAPA{Ejemplo de tabla elaborada para mostrar el formato. En APA se evitan las líneas verticales y los bordes alrededor de cada celda.}
\end{table}

\section{Concepto 4}

La Figura~\ref{fig:ejemplo} muestra un ejemplo de figura con número, título y nota.

\begin{figure}[htbp]
  \caption{Ejemplo de una figura en formato APA}
  \label{fig:ejemplo}
  \begin{center}
    % Sustituye el recuadro por tu imagen real:
    % \includegraphics[width=0.80\textwidth]{imagenes/mi_imagen.png}
    \fbox{\rule{0pt}{5cm}\rule{0.80\textwidth}{0pt}}
  \end{center}
  \notaAPA{Reemplaza este texto por la explicación de la figura y, si corresponde, por la atribución de la fuente.}
\end{figure}

\section{Clasificación de ...}

Introduce la clasificación antes de crear subsecciones. Explica el criterio usado para distinguir las clases.

\subsection{Clase 1}

Explica aquí la primera clase.

\subsection{Clase 2}

Explica aquí la segunda clase.

\section{Estado del arte}

Puedes cerrar el capítulo con un estado del arte o una síntesis. Conviene comparar enfoques, resultados, limitaciones y vacíos de los trabajos previos.

% Ejemplo de cita textual de 40 palabras o más:
% \begin{citaAPA}
% Aquí se coloca una cita textual extensa de cuarenta palabras o más. No se usan
% comillas y todo el bloque queda sangrado 1.27 cm desde el margen izquierdo.
% Al final se añade la fuente. \parencite[100]{ejemplo2020}
% \end{citaAPA}
